\label{ch:4}

\section{Design Process or Methodology Overview}
\label{sec:4.1}

This chapter explains the design of the system and why it takes this form. The target is a reply that makes a person feel understood, produced on a single 8 GB consumer GPU, with the emotional decision recorded rather than buried in a model. Three guidelines shaped every choice: \textbf{recognise the feeling before writing anything}; \textbf{give the generator one instruction at a time}, because a model told to answer several feelings at once answers none of them well; and \textbf{adopt nothing the measurements do not support}. The work ran in two phases, shown in Figure~\ref{fig:4.1}. The first compared five detector architectures on the text the system actually receives. The second built the response side on that detector and compared generators, four ways of passing the detector's output on, and monolithic baselines that reply without any detector.

\begin{figure}[H]
\centering
\includegraphics[width=0.95\linewidth]{ch4_fig1.png}
\caption{The evidence-gated design process followed in this work}
\label{fig:4.1}
\end{figure}

\subsection{The Proposed System}

The final system is a two-stage pipeline built from two models (Figure~\ref{fig:4.2}). Every part has a purpose.

\begin{itemize}
  \item \textbf{d1, the emotion detector:} a RoBERTa-base encoder fine-tuned on the 32 emotions of EmpatheticDialogues \cite{rashkin2019} reads the user's message, or the last four user turns in conversation, and outputs a probability for every emotion; only its three most probable emotions leave this stage. An encoder is used rather than a prompted generator because fine-tuned encoders outperform prompted models at emotion recognition \cite{belay2025}, and 125 million parameters cost almost nothing beside a 3.8-billion-parameter generator.
  \item \textbf{Attention cues}: the words the detector attended to, taken from its last attention layer in the same forward pass, so the reply can refer to what the person wrote and not only to an emotion word.
  \item \textbf{Stage 1, emotion choice:} in a short deterministic call, Phi-4-mini-instruct \cite{abouelenin2025phi4mini} picks the one emotion among d1's three candidates that best fits the message, resolving the detector's uncertainty before anything is written.
  \item \textbf{Prompt builder:} the chosen emotion selects one of 32 pieces of guidance, combined with the cue words and the three response rules into a single system prompt.
  \item \textbf{Stage 2, reply generation:} the same generator, in 4-bit form, writes the reply, reading the last ten exchanges.
\end{itemize}
Every intermediate decision is recorded: the three candidates with their confidences, the emotion finally used, and the cue words. That record is what a single end-to-end model cannot produce.

\begin{figure}[H]
\centering
\includegraphics[width=0.75\linewidth]{ch4_fig2.png}
\caption{The complete two-stage pipeline of the proposed system}
\label{fig:4.2}
\end{figure}

\subsection{Design Space Explored}

Because this thesis asks whether that separation is merited, the design space deliberately includes designs that would make the detector unnecessary. Four families were built and run under identical conditions:

\textbf{Detector alternatives:} one decoder and four encoders trained on the same data with the same settings (Section 4.2.1).

\textbf{Generator alternatives: }three instruction-tuned generators of 1.5 to 4 billion parameters, with the detector, prompt and conversations fixed (Section 4.2.2).

\textbf{Integration Alternatives:} four ways of feeding d1's output into d2, plus the two-stage design finally adopted (Section 4.2.3).

\textbf{Monolithic and retrieval alternatives:} five designs that reply without a detector, and one that retrieves human examples instead (Section 4.2.5).

All ran on the same 300 held-out conversations with the same decoding settings, so differences come from the design and not from the data. Chapter 5 reports how each performed.

\subsection{Evaluation Strategy}

Since no automatic measure captures empathy, every design is evaluated on five independent signals: rule adherence (an automatic check of whether a reply acknowledges the feeling, asks exactly one question and stays within 120 words); the WASSA 2024 empathy classifier; emotion validation, whether the emotion expressed in the reply matches the conversation's label; two AI judges from unrelated developers, scoring blind; and a study with human raters. Detectors are judged on top-1 accuracy, top-5 accuracy and macro-F1 over the full test partition. Every system answers the same conversations, so comparisons are paired: McNemar's exact test for binary outcomes \cite{mcnemar1947}, Wilcoxon signed-rank for scores \cite{wilcoxon1945}, bootstrap intervals with 10,000 resamples \cite{efron1993} and Holm correction across families \cite{holm1979}. A difference counts as significant only when the corrected p-value is below 0.05 and the interval excludes zero.

\subsubsection{Selection of the AI-Judge Panel}

Language models used as judges agree with human raters more than 80\% of the time \cite{zheng2023judge}, but they favour their own outputs \cite{panickssery2024}, so the panel was designed against that bias: \textbf{Nemotron 3 Ultra} (NVIDIA) and \textbf{gpt-oss-120b }(OpenAI) share no family with our generator (Microsoft) or with the Qwen baselines (Alibaba). The reported score averages the two, since a panel drawn from different families carries less bias than a single large judge \cite{verga2024}. Each scores empathy, relevance and safety from 1 to 10 against written anchors at temperature 0, seeing only the user's message and one reply, never the system's name (Table~\ref{tab:4.1}).

\begin{table}[H]
\centering
\scriptsize
\begin{tabularx}{\textwidth}{@{}lXXXXX@{}}
\toprule
\textbf{Judge} & \textbf{Developer} & \textbf{Size} & \textbf{Access} & \textbf{Role} & \textbf{Why selected} \\
\midrule
Nemotron 3 Ultra & NVIDIA & 550B total / 55B active & OpenRouter free tier & Every judged system & Flagship scale; unrelated to either family \\
gpt-oss-120b & OpenAI (open weights) & About 120B & Groq free tier & Every judged system & Open weights make it reproducible \\
\bottomrule
\end{tabularx}
\caption{The AI-judge panel}
\label{tab:4.1}
\end{table}

\subsubsection{Use of the WASSA Empathy Classifier}

Independently of the rule checks and the judges, replies are scored by the hyy33 system from the WASSA 2024 shared task \cite{yang2024hyy33}. It suits this project for three reasons: it was trained on third-party human empathy ratings of single conversation turns \cite{giorgi2024wassa}, the same unit as our replies; it is peer-reviewed and publicly released, correlating with human ratings at r = 0.544 on that task's test set; and it comes from a different group, dataset and model family from every part of our system. Its limits are equally clear: it was not the best system in its own task, it was trained on news articles rather than supportive replies, and it produces coarse classes. Section 5.4.4 shows it reproduces the judges' ranking well but is too noisy per reply, so it is used only to compare system averages.

\section{Preliminary Design or Design (Model) Specification}
\label{sec:4.2}

This section describes the candidates for each component and the full specification of the one finally chosen.

\subsection{Emotion Detector Candidates (d1)}

Five detectors were trained on the same data with the same settings, varying only the architecture (Table~\ref{tab:4.2}).

\textbf{GPT-2 Medium (the initial detector):} a 355-million-parameter decoder-only transformer with a classification head on its last hidden state, trained on the \textit{situation description}, the paragraph a crowdworker wrote before the conversation began. Two properties make it a weak fit: a left-to-right decoder sees each token only in the context of those before it, whereas classification benefits from reading the whole message at once; and the situation description names the feeling more explicitly than a person's own message, so a detector trained on it learns an easier task than the one it meets in use.

\textbf{BERT-base }and \textbf{ELECTRA-base: }110-million-parameter bidirectional encoders, the first pre-trained with masked-language modelling \cite{devlin2019} and the second by replaced-token detection, which trains on all positions rather than the 15\% masking covers.

\textbf{RoBERTa-base:} a 125-million-parameter encoder, otherwise identical to BERT but trained far longer on about 160 GB of text with dynamic masking and no next-sentence objective \cite{liu2019roberta}. It is the size the memory budget favours, leaving nearly all of the 8 GB for the generator.

\textbf{RoBERTa-large: }the same model at 355 million parameters, included for a cost question rather than an architectural one.

All five use a 32-way classification head with cross-entropy loss and the settings of Section 4.4. Selection used top-1 and top-5 accuracy on the 2,542-conversation test partition, weighting top-5 heavily because the design consumes a ranking rather than a single guess.

\begin{table}[H]
\centering
\scriptsize
\begin{tabularx}{\textwidth}{@{}lXXXX@{}}
\toprule
\textbf{Candidate} & \textbf{Architecture} & \textbf{Parameters} & \textbf{Input it was trained on} & \textbf{Why it was included} \\
\midrule
GPT-2 Medium & Decoder-only transformer & 355M & Situation description & The initial pipeline's detector \\
BERT-base & Bidirectional encoder & 110M & First user message & The original masked-language encoder \\
ELECTRA-base & Bidirectional encoder & 110M & First user message & Replaced-token detection, often stronger than BERT at this size \\
RoBERTa-base & Bidirectional encoder & 125M & First user message & Longer pre-training on 160 GB; the size the memory budget favours \\
RoBERTa-large & Bidirectional encoder & 355M & First user message & Tests whether a larger detector is worth its cost \\
\bottomrule
\end{tabularx}
\caption{Detector candidates, their architectures and why each was included}
\label{tab:4.2}
\end{table}

\subsection{Response Generator Candidates (d2)}

Only the generator changed; the detector, prompt, rules, conversations and decoding settings were fixed (Table~\ref{tab:4.3}).

\textbf{Qwen2.5-1.5B-Instruct} is the initial pipeline's generator \cite{qwen2024}, the lower bound on size. \textbf{Gemma-3-4B-it} tests whether capacity alone produces supportive replies, though with no system role in its chat template instructions must be folded into the user turn. \textbf{Phi-4-mini-instruct} is trained with a strong emphasis on instruction following \cite{abouelenin2025phi4mini}, which three explicit rules and a word limit demand, and is loaded with 4-bit NF4 quantisation \cite{dettmers2023qlora} to fit beside the detector. \textbf{Phi-2} was dropped before evaluation: with no chat template the rules could not be enforced.

\begin{table}[H]
\centering
\scriptsize
\begin{tabularx}{\textwidth}{@{}lXXXX@{}}
\toprule
\textbf{Candidate} & \textbf{Parameters} & \textbf{Loading} & \textbf{Chat template} & \textbf{Why it was included} \\
\midrule
Qwen2.5-1.5B-Instruct & 1.5B & 16-bit & Yes & The initial pipeline's generator; the size a small GPU runs \\
Gemma-3-4B-it & 4B & 4-bit NF4 & Yes, but no system role & A different family; tests whether size alone helps \\
Phi-4-mini-instruct & 3.8B & 4-bit NF4 & Yes & Emphasis on instruction following, which the rules demand \\
Phi-2 & 2.7B & 16-bit & No & Dropped: without a chat template it cannot be steered \\
\bottomrule
\end{tabularx}
\caption{Generator candidates, with the detector, prompt and conversations held fixed}
\label{tab:4.3}
\end{table}

\subsection{Integration Design: How d1's Output Reaches d2}
\label{sec:4.2.3}

The detector produces a full probability distribution, so how that output reaches the generator is itself a design choice (Table~\ref{tab:4.4}). Four variants were built on the first pipeline: \textbf{argmax}, passing only the most probable emotion; \textbf{confidence-gated}, the same label with tentative wording when the detector is unsure; \textbf{distribution}, every emotion above probability 0.05; and \textbf{candidate-pick}, the top three with the generator choosing one and replying in a single call.

These were compared in Section 5.2.4, which led to the\textbf{ two-stage} design used in the final system. Two facts motivate it: the correct emotion sits among the detector's top three far more often than it is the single most likely, so a larger candidate set carries real information; and several emotions in one prompt give conflicting advice, since answering disappointment well is not the same as answering pride well. A short, greedy call removes the uncertainty first and leaves the generation prompt describing exactly one feeling.

\begin{table}[H]
\centering
\footnotesize
\begin{tabularx}{\textwidth}{@{}lXXX@{}}
\toprule
\textbf{Variant} & \textbf{What reaches the generator} & \textbf{Design rationale} & \textbf{Risk it carries} \\
\midrule
Argmax & The single most probable emotion & The simplest link; one emotion, one coherent instruction & A wrong top-1 label misdirects the whole reply \\
Confidence-gated & The top emotion, hedged when d1 is unsure & Lets the reply hedge instead of asserting a doubtful feeling & Depends on the generator actually hedging when told to \\
Distribution & Every emotion above probability 0.05 & Passes the detector's uncertainty on in full & Several emotions at once give contradictory guidance \\
Candidate-pick & The top three, resolved inside the generation call & One call, so no extra latency & The model must resolve the emotion and obey the rules at once \\
Two-stage (selected) & The top three, resolved in a separate short call & A separate call keeps the generation prompt coherent & One extra model call per reply, at 0.31 s \\
\bottomrule
\end{tabularx}
\caption{Ways of passing the detector's output to the generator}
\label{tab:4.4}
\end{table}

\subsection{Attention-Based Cue Extraction}

So that the reply addresses what the person wrote and not merely an emotion word, the design copies out the words the detector itself used (Figure~\ref{fig:4.3}). Let \(A \in \mathbb{R}^{H \times T \times T}\) be the encoder's last-layer attention, with \(H\) heads and \(T\) tokens. The cue score of token \(t\) is the attention the classification token pays to it, averaged over heads:

\[
s(t) = \frac{1}{H}\sum_{h=1}^{H} A_{h,0,t}.
\]

Sub-word pieces are merged back into words by adding their scores, special tokens are dropped, and the highest-scoring 40\% of words are kept, with a floor of three and a ceiling of fourteen so that neither a very short nor a very long message produces a useless cue list. The result is a few phrases such as "moved away" or "my dog", which the prompt builder places in the instruction.

One implementation detail is essential: the attention implementation must be requested explicitly, because the library's default returns no weights and raises no error, silently falling back to text splitting. The system therefore checks at start-up whether weights are returned and stops if they are not.

\begin{figure}[H]
\centering
\includegraphics[width=0.95\linewidth]{ch4_fig3.png}
\caption{Extraction of cue words from the detector's final-layer attention}
\label{fig:4.3}
\end{figure}

\subsection{Monolithic and Retrieval Alternatives}

Six designs were built without any separate detector (Table~\ref{tab:4.5}) to test whether one is needed. Three prompt a single model to detect the emotion silently and then reply under the same three rules, at 1.5 billion parameters, at 7.6 billion in 4-bit, and with our own 3.8-billion-parameter generator; the last is the most telling, being the proposed system without the detector. A fourth asks the generator to be generically caring with no emotion at all. A fifth fine-tunes Qwen2.5-1.5B with LoRA \cite{hu2022lora} on pairs of user messages and human replies. The sixth is retrieval-augmented \cite{lewis2020rag}: one or two training conversations with the same detected emotion are retrieved using all-MiniLM-L6-v2 sentence embeddings \cite{reimers2019sbert} and shown to the generator as examples.

\begin{table}[H]
\centering
\footnotesize
\begin{tabularx}{\textwidth}{@{}lXX@{}}
\toprule
\textbf{Baseline} & \textbf{How it is built} & \textbf{What it is there to test} \\
\midrule
Zero-shot generator & Qwen2.5-1.5B, generic caring prompt, no detector & Whether a detector adds anything over plain prompting \\
Fine-tuned monolithic & Qwen2.5-1.5B fine-tuned with LoRA on message and reply pairs & Whether learning the dataset's replies beats prompting \\
Single model, same size & Qwen2.5-1.5B told to identify the emotion silently, then reply & The core question: decoupled against monolithic at equal size \\
Single model, larger & Qwen2.5-7B in 4-bit with the same prompt & What is given up by staying inside the 8 GB budget \\
Single model, our generator & Phi-4-mini alone with the same prompt and no detector & Whether the gains come from the design or the generator \\
Retrieval-augmented & Similar conversations retrieved with all-MiniLM-L6-v2, shown as examples & Whether human examples improve the reply \\
\bottomrule
\end{tabularx}
\caption{Designs built without a separate detector, and what each one tests}
\label{tab:4.5}
\end{table}

\subsection{Specification of the Selected Design}

\textbf{Detector:} RoBERTa-base with a 32-way classification head, fine-tuned on first messages with the settings of Table~\ref{tab:4.8}. It reads one message when answering a single message, and the user's last four turns joined together in conversation; Section 5.3.2 gives the measurements behind that window.

\textbf{Stage 1, emotion choice: }Phi-4-mini receives the message and the three candidates and is asked:\textit{ "An automatic detector narrowed the person's likely feeling down to these options: \{a\}, \{b\}, \{c\}. Read their message and decide which ONE of those options best fits what they are actually expressing. Reply with ONLY that single word."} Decoding is greedy and capped at 10 new tokens. If the answer names none of the three the system falls back to the detector's first choice, though in practice every answer was valid.

\textbf{Stage 2, reply generation:} the system prompt is assembled in a fixed order from the chosen emotion, the cue words, the guidance for that emotion, and the three universal rules (Figure~\ref{fig:4.4}). For \textit{lonely}, the guidance tells the model to "acknowledge the isolation itself, gently" and not to "offer a surface-level fix like 'go meet people'"; the rules then require it to validate the feeling in two or three sentences, offer a reassuring thought specific to it, and end with exactly one short question, all within 100 words. Guidance of this kind exists for all 32 emotions, and Table~\ref{tab:4.9} gives the decoding settings.

\begin{figure}[H]
\centering
\includegraphics[width=0.95\linewidth]{ch4_fig4.png}
\caption{Assembly of the stage-2 system prompt from four inputs}
\label{fig:4.4}
\end{figure}

Figure~\ref{fig:4.5} prints the prompt as the generator receives it for a message read as \textit{lonely}.

\begin{figure}[H]
\centering
\includegraphics[width=0.95\linewidth]{ch4_fig5.png}
\caption{The system prompt as the generator actually receives it}
\label{fig:4.5}
\end{figure}

\subsection{The Conversational Agent}

The evaluation in Chapter 5 answers one message at a time, but the system people use holds a conversation. The agent is the same two-stage pipeline run once per turn, with four additions the single-reply system does not need (Table~\ref{tab:4.6}).

\textbf{Two memories for two different models: }the detector is a classifier trained on single short messages, so it reads a bounded window, the user's last four turns, and only the user's own text. The generator has a large context window, so it receives the last ten exchanges verbatim; summarising them would compress away the emotional detail the reply depends on.

\textbf{Rules that change after the first turn:} the three rules are unchanged on the first turn, keeping it comparable with every single-turn result in Chapter 5. From the second turn the question is no longer compulsory, and the generator is asked not to repeat an acknowledgement or reuse earlier phrasing.

\textbf{Emotions that move and an audit trail: }each turn's emotion is compared with the previous turn's, and when it changes the agent names the change and answers what the person feels now. Every turn records the three candidates with their probabilities, the emotion chosen, the cue words, any shift, and the reply. An end-to-end model can produce the same reply but cannot show which feeling it answered, nor allow any intervention before it.

\begin{table}[htbp]
\centering
\footnotesize
\begin{tabularx}{\textwidth}{@{}lXX@{}}
\toprule
\textbf{Part of the agent} & \textbf{Setting} & \textbf{Why this setting} \\
\midrule
Detector input & Last 4 user turns, joined, user text only & 4 was the widest window that degraded no scenario (Section 5.3.2) \\
Attention cues & Current message only & Cues should point at what was just said, not at turns already passed \\
Generator memory & Last 10 exchanges, verbatim & Support conversations are short, and summarising loses emotional detail \\
First-turn rules & The three single-turn rules, unchanged & Keeps the first exchange comparable with Chapter 5 \\
Later-turn rules & At most one question, answer what is new, do not reuse phrasing & A question every turn turns the conversation into an interrogation \\
Emotion shift & Previous emotion recorded; the prompt names the change & Feelings move, and replying to the earlier one reads as not listening \\
Audit trail & Per turn: candidates and probabilities, chosen emotion, cue words, shift, reply & The inspectability the design exists for \\
\bottomrule
\end{tabularx}
\caption{The conversational agent: what each setting is and why}
\label{tab:4.6}
\end{table}

\section{Data Collection and Preparation}
\label{sec:4.3}

\subsection{Source of the Data}

All data comes from EmpatheticDialogues \cite{rashkin2019}, a corpus in which one crowdworker describes a situation tied to one of 32 emotions and another replies as a listener. The released files hold 23,144 conversations and 88,687 utterances, already split so that every conversation from one speaker stays in a single set, preventing a speaker's style leaking from training into test. The corpus is used under its non-commercial research licence \cite{ed_github}. Figure~\ref{fig:4.6} shows the route from those files to the four datasets this work uses.

\begin{figure}[htbp]
\centering
\includegraphics[width=0.95\linewidth]{ch4_fig6.png}
\caption{From the released corpus to the datasets used in this work}
\label{fig:4.6}
\end{figure}

\textbf{Why not GoEmotions: }the project began on GoEmotions \cite{demszky2020}: 54,263 Reddit comments across 27 emotions and a neutral class, processed into a 28-label multi-label set and rewritten as dialogue with templated replies. It was abandoned for four reasons, each measured on that prepared data rather than assumed.

\begin{itemize}
\item \textbf{ It has no replies:} a comment carries only a label, so there is nothing to train or evaluate a generator against; the replies we wrote for it were our own.

\item \textbf{ A third of it has no emotion to answer:} \textit{neutral }is the largest class at 32.75\% of the corpus, some 18,000 comments.

\item \textbf{ The emotions it does carry are the wrong ones:} its frequent classes are reactions to online comments (\textit{admiration }9.4\%, \textit{approval }6.8\%), while the states a support conversation must answer are its rarest: \textit{grief }appears 96 times against \textit{admiration}'s 5,122, and it has no \textit{lonely}, \textit{devastated}, \textit{terrified }or \textit{apprehensive }at all.

\item \textbf{ Its labels overlap: }16.25\% of examples carry more than one emotion, most often \textit{anger }and \textit{annoyance}, which does not fit a design that conditions one reply on one emotion.
\end{itemize}

EmpatheticDialogues answers all four: every conversation is a first-person disclosure written to be replied to, it includes the listener's reply, and it carries exactly one of 32 emotions chosen by the speaker beforehand. The change of corpus is why this thesis uses 32 labels rather than the 28 originally planned.

\subsection{Data Cleaning}

Three cleaning steps each address a specific defect in the released files. The literal string \_comma\_ is restored to a comma, both in the text a model reads and in the text a rater sees. A row that does not parse as valid CSV, a few because of unescaped quotation marks, is dropped rather than repaired. Any row with an empty utterance or label is removed, and text and labels are trimmed of surrounding whitespace.

\subsection{Data Reduction}

The detector reads one message, not a whole conversation, so the first utterance of each conversation is kept with that conversation's label and the rest discarded. This is the most consequential preparation decision in the project, replacing the initial pipeline's choice of training on the \textit{situation description}, which is written after the event and names the feeling almost explicitly; Section 5.2.2 measures what the mismatch cost. Figure~\ref{fig:4.7} shows one conversation before and after preparation.

A second reduction defines the evaluation set: 300 conversations drawn from the test partition with seed 42, used by every generation experiment in this thesis. Of those, 64 are selected, the two shortest and two longest messages for each emotion, so that all emotions are represented within the judges' daily request limits.

\begin{figure}[H]
\centering
\includegraphics[width=0.95\linewidth]{ch4_fig7.png}
\caption{One conversation before and after preparation}
\label{fig:4.7}
\end{figure}

\subsection{Data Transformation}

Text is transformed only as far as each model requires. Messages for the detector are tokenised with the model's own tokeniser and truncated to 128 tokens, which covers 99.7\% of first messages; the 32 emotion names are mapped to integer labels in a fixed order so predictions are comparable across candidates. Except for Gemma, the user's message is inserted into the model's chat template with the system prompt of Section 4.2.6 in the system role. No lower-casing, stemming or stop-word removal is applied anywhere, because those steps would strip the very cues, capitalisation and punctuation among them, that signal feeling.

\subsection{Data Integration}

For each evaluation conversation the user's first message is paired with the human listener's reply, the latter kept only as a reference for reading and never as a target the systems are scored against. Generated replies are re-joined to their conversation by identifier, so the reply, the gold emotion, the detector's prediction, the rule checks and the judge scores sit in one row. That structure is what makes the paired tests of Chapter 5 possible.

\subsection{Summary of the Prepared Data}

Table~\ref{tab:4.7} shows what preparation produced: 17,839 labelled messages to train the detector, 2,763 to select the checkpoint, 2,542 to measure it, 300 conversations answered by every system, 64 of those scored by both AI judges, and 8 placed before the human raters.

\begin{table}[htbp]
\centering
\footnotesize
\begin{tabularx}{\textwidth}{@{}lXXX@{}}
\toprule
\textbf{Dataset produced} & \textbf{Size} & \textbf{Unit} & \textbf{Used for} \\
\midrule
Detector training set & 17,839 & First message with its label & Fine-tuning every detector candidate \\
Detector validation set & 2,763 & First message with its label & Early stopping and checkpoint selection \\
Detector test set & 2,542 & First message with its label & Reported detector accuracy \\
Generator training pairs & 17,749 & First message with the human reply & The LoRA fine-tuned baseline only \\
Reply evaluation set & 300 & Test conversation & Every generation experiment, identical across systems \\
AI-judge sample & 64 & Test conversation & Blind scoring by both judges, 2 per emotion \\
Human-study sample & 8 & Test conversation & The rating study of Chapter 5 \\
\bottomrule
\end{tabularx}
\caption{The datasets produced by preparation and what each is used for}
\label{tab:4.7}
\end{table}

\section{Implementation of the Selected Design}
\label{sec:4.4}

\subsection{Detector Training}

Every detector candidate was trained end to end with cross-entropy loss over the 32 emotions using the setup of Table~\ref{tab:4.8}. Training runs for at most 4 epochs with early stopping after one epoch without improvement, and the reported checkpoint is the best rather than the last, so a run that begins to over-fit is not carried into the results. The learning rate warms up linearly over the first 6\% of steps; this is RoBERTa's own fine-tuning procedure and is not cosmetic, since without it the RoBERTa-large run diverged to 5.0\% accuracy until the cause was found (Section 3.7). Mixed 16-bit precision and gradient accumulation let an effective batch of 16 fit in 8 GB. All five candidates used the same script with only the model name changed, and seed 42 throughout, so the comparison is between architectures and not training recipes.

\begin{table}[H]
\centering
\footnotesize
\begin{tabularx}{\textwidth}{@{}lXX@{}}
\toprule
\textbf{Setting} & \textbf{Value} & \textbf{Reason} \\
\midrule
Backbone & RoBERTa-base, 32-way head & Selected in Section 5.2.2 \\
Training input & First user message of each conversation & The text the detector receives in use \\
Maximum input length & 128 tokens & Covers 99.7\% of first messages (mean 46.5) \\
Learning rate & 2e-5 & Inside the range recommended for BERT-family encoders \\
Effective batch size & 16 (8 per step x 2 accumulation) & Largest that fits in 8 GB at this length \\
Epochs & Up to 4, early stopping with patience 1 & Best checkpoint chosen by validation accuracy \\
Warm-up & Linear over the first 6\% of steps & RoBERTa's own setting; without it a large run diverged \\
Weight decay / gradient clipping & 0.01 / 1.0 & Standard regularisation for this family \\
Precision / seed & 16-bit mixed precision / 42 & Fits the memory budget and keeps runs reproducible \\
\bottomrule
\end{tabularx}
\caption{Detector training configuration}
\label{tab:4.8}
\end{table}

\subsection{Generator Configuration}

The generator is never trained in the final system: it is used zero-shot and steered only by its system prompt. Two implementation decisions make that fit the memory budget. The model is loaded with 4-bit NF4 quantisation and double quantisation with 16-bit compute \cite{dettmers2023qlora}, which cuts memory roughly fourfold while keeping the quality of 16-bit weights. Decoding for every system in this thesis samples at temperature 0.7 with top-p 0.9 rather than using greedy or beam search, which are known to produce repetitive text \cite{holtzman2020}. These settings are identical everywhere, so any difference between systems comes from the design and not from decoding.

\begin{table}[H]
\centering
\footnotesize
\begin{tabularx}{\textwidth}{@{}lXX@{}}
\toprule
\textbf{Setting} & \textbf{Value} & \textbf{Reason} \\
\midrule
Generator & Phi-4-mini-instruct (3.8B) & Selected in Section 5.2.3 \\
Quantisation & 4-bit NF4, double quantisation, 16-bit compute & Keeps 16-bit quality at about a quarter of the memory; peak 3.40 GB \\
Stage-1 decoding & Greedy, at most 10 new tokens & The emotion choice must be deterministic and short \\
Stage-2 temperature & 0.7 & The developer's default for this family \\
Stage-2 top-p & 0.9 & Nucleus sampling avoids the repetition of greedy decoding \\
Maximum new tokens & 250 & Well above the 100-word request; longest reply was 119 tokens \\
Candidates passed to stage 1 & Top 3 emotions & The true emotion is among the top three far more often than first \\
Attention cues & Final-layer [CLS] attention, top 40\% of words & Gives the generator the words the detector used \\
Seed & 42 & Identical decoding for every system \\
\bottomrule
\end{tabularx}
\caption{Generation configuration, identical for every system compared}
\label{tab:4.9}
\end{table}

\subsection{Implementing the Two Stages}

Stage 1 requests at most 10 tokens and checks the answer against the three candidate names, falling back to the detector's first guess if none matches. In stage 2 the system prompt is assembled in the order shown in Figure~\ref{fig:4.4}, and the sampling call allows 250 new tokens so that no reply is cut off mid-sentence; the longest the final system produced was 119 tokens. Three model-specific details must be handled explicitly or they silently corrupt results: Phi-4-mini's end-of-turn token differs from its end-of-text token, so generation must stop on the former; Gemma has no system role in its chat template; and Phi-4-mini merges its attention projections into a single layer, which matters when naming the modules for low-rank adaptation.

\subsection{Implementing the Baselines}

The monolithic baselines share the pipeline's scaffolding, the same 300 conversations, decoding settings, rule checks and output format, with the detector removed and the prompt asking the model to work out the feeling itself. The LoRA baseline is the one design that trains a generator, using rank 16 and alpha 32 on the attention projections with a learning rate of 2e-4 \cite{hu2022lora}. One correction is worth recording: the padding token must be the model's end-of-text token, because using the end-of-turn token removes it from the loss, and a model that never sees it learns not to stop. The first version had that defect and was retrained once the cause was found (Section 5.3.3).

\subsection{Complete Parameter Listing}

Every parameter was fixed before the final comparison, justified by published practice, a measurement on our own data, or a hardware limit. Table~\ref{tab:4.10} gives the values needed to reproduce this work. No tuning was carried out on the test set.

\begin{table}[H]
\centering
\footnotesize
\renewcommand{\arraystretch}{1.1}
\begin{tabularx}{\textwidth}{@{}>{\raggedright\arraybackslash}p{0.23\textwidth}>{\raggedright\arraybackslash}p{0.22\textwidth}>{\raggedright\arraybackslash}X@{}}
\toprule
\textbf{Parameter} & \textbf{Value} & \textbf{Why this value} \\
\midrule
Detector: learning rate & 2e-5 & Inside the ranges recommended for BERT (2e-5 to 5e-5) and used for RoBERTa (1e-5 to 3e-5) \\
Detector: effective batch size & 16 (8 per step x 2 accumulation) & The batch size recommended for BERT/RoBERTa fine-tuning; 8 per step fits in 8 GB \\
Detector: epochs and checkpoint choice & Up to 4, early stopping (patience 1), best by validation accuracy & BERT recommends 2-4 epochs; keeping the best checkpoint prevents over-fitting \\
Detector: precision and seed & fp16 mixed precision / 42 & Fits the 8 GB GPU and trains faster; fixed seed for reproducibility \\
Detector: RoBERTa-large learning rate and batch & 1e-5 / 4 & The base settings diverged; 1e-5 is the lowest of RoBERTa's range and batch 4 fits in memory \\
Attention cues: how much is kept & Top 40\% of content tokens, floor 3, ceiling 14, at most 3 phrases & A fixed count returns nearly the whole message for short inputs, which adds nothing; the cap keeps the prompt short \\
Generator: quantisation & 4-bit NF4, double quantisation, fp16 compute & NF4 keeps 16-bit quality at about a quarter of the memory; measured peak 3.40 GB \\
\bottomrule
\end{tabularx}
\caption{All parameters list with reasoning}
\label{tab:4.10}
\end{table}

\begin{table}[H]
\centering
\footnotesize
\renewcommand{\arraystretch}{1.1}
\begin{tabularx}{\textwidth}{@{}>{\raggedright\arraybackslash}p{0.23\textwidth}>{\raggedright\arraybackslash}p{0.22\textwidth}>{\raggedright\arraybackslash}X@{}}
\toprule
\textbf{Parameter} & \textbf{Value} & \textbf{Why this value} \\
\midrule
Generator: temperature & 0.7 & Qwen2.5-Instruct's published default, then held fixed for every system and generator so differences come from the design, not decoding \\
Agent: how much the detector reads & Last 4 user turns & Best or joint best in 3 of 4 long scenarios; 6 turns degraded 2 of them (Section 5.3.2) \\
Agent: how much the generator remembers & Last 10 exchanges, verbatim & Support conversations are short; verbatim history avoids losing emotional detail \\
LoRA baseline: rank and alpha & 16 / 32 & Low ranks suffice to adapt attention weights; alpha = 2 x rank scales the update \\
LoRA baseline: learning rate, dropout, weight decay & 2e-4 / 0.1 / 0.01 & LoRA's own settings; QLoRA uses 2e-4 and dropout 0.1 for models up to 13B \\
LoRA baseline: epochs and checkpoint choice & Up to 3, best epoch by validation loss, early stopping (patience 1) & Validation loss rose after epoch 1 (1.4607 to 1.4626), so epoch 1 was kept \\
\bottomrule
\end{tabularx}
\caption*{Table 4.10: All parameters list with reasoning (continued)}
\end{table}

\subsection{Verification and Safeguards}

A small trial run preceded the full evaluation, checking that every design produced well-formed output and that the model-specific details above were handled. Two automatic safeguards then protect the chosen design: if the detector returns no attention weights the cue mechanism stops rather than failing silently, and every stage-1 choice is written to the output file, so a malformed choice is noticed rather than swallowed by the fallback. The repository also holds the scored outputs of all 4,200 generated replies, so every table in Chapter 5 can be rebuilt by script.
